% ============================================================
%  Rapport minimaliste — TP RCR1
%  Master 1 SII · USTHB · 2025-2026
%
%  Version condensée : pour chaque TP, on décrit le domaine
%  choisi, la représentation des connaissances, puis l'exploitation.
% ============================================================
\documentclass[11pt, a4paper]{article}

\usepackage[utf8]{inputenc}
\usepackage[T1]{fontenc}
\usepackage[french]{babel}
\usepackage{lmodern}
\usepackage[top=2.2cm, bottom=2.2cm, left=2.5cm, right=2.5cm]{geometry}
\usepackage{amsmath, amssymb}
\usepackage{booktabs}
\usepackage{tabularx}
\usepackage{xcolor}
\usepackage{listings}
\usepackage[most]{tcolorbox}
\usepackage[colorlinks=true, linkcolor=blue!60!black, urlcolor=blue!60!black,
            citecolor=blue!60!black]{hyperref}

% Couleurs
\definecolor{usthbblue}{RGB}{0,63,130}
\definecolor{codebg}{RGB}{248,248,248}
\definecolor{codegreen}{RGB}{40,140,40}
\definecolor{codeblue}{RGB}{0,90,180}

% Styles tcolorbox
\tcbset{
  defbox/.style={
    colback=usthbblue!5, colframe=usthbblue!60,
    fonttitle=\bfseries, boxrule=0.5pt, arc=4pt,
  },
}

\lstset{
  backgroundcolor=\color{codebg},
  basicstyle=\ttfamily\footnotesize,
  commentstyle=\color{codegreen}\itshape,
  keywordstyle=\color{codeblue}\bfseries,
  breaklines=true,
  frame=single, rulecolor=\color{gray!40},
  xleftmargin=8pt, xrightmargin=4pt,
  showstringspaces=false,
  tabsize=2, language=Java,
}

\newcommand{\code}[1]{\texttt{#1}}

% Symboles
\newcommand{\nec}{\Box}
\newcommand{\pos}{\Diamond}
\newcommand{\sq}{\sqsubseteq}

% ============================================================
\begin{document}

% ============================================================
% PAGE DE GARDE
% ============================================================
\begin{titlepage}
\begin{center}

% En-tête université
{\large\bfseries République Algérienne Démocratique et Populaire}\\[0.3em]
{\large Ministère de l'Enseignement Supérieur et de la Recherche Scientifique}\\[1.5em]

\rule{\linewidth}{0.6pt}\\[0.5em]
{\large\bfseries Université des Sciences et de la Technologie Houari Boumediene}\\
{\normalsize Faculté d'Informatique · Département SI}\\[0.3em]
\rule{\linewidth}{0.6pt}

\vspace{2em}

% Discipline
{\color{usthbblue}\large Master 1 Systèmes d'Information et d'Intelligence Artificielle}\\[0.5em]
{\normalsize Module : \textbf{RCR1 --- Représentation des Connaissances et Raisonnement}}

\vspace{3em}

% Titre principal
\begin{tcolorbox}[defbox, width=0.92\linewidth, halign=center]
  {\LARGE\bfseries Implémentation de Formalismes de}\\[0.4em]
  {\LARGE\bfseries Représentation des Connaissances}\\[0.4em]
  {\LARGE\bfseries et Raisonnement}\\[1em]
  {\large Logique Modale · Logique des Défauts · Logique Classique}\\
  {\large Logique de Description · Réseaux Sémantiques}
\end{tcolorbox}

\vspace{3em}

% Domaines d'application
{\normalsize Domaines d'application :
  \textbf{Industrie musicale algérienne} et \textbf{Système judiciaire}}

\vspace{3em}

% Auteurs
\begin{tabular}{rl}
  \textbf{Réalisé par :} & D. Arslan \\
                          & [Prénom Nom Binôme 2] \\[1em]
  \textbf{Encadré par :} & [Nom de l'encadrant] \\[1em]
  \textbf{Groupe :}      & [Numéro du groupe]
\end{tabular}

\vfill

% Outils
{\small\color{gray}
  TweetyProject 1.27 · OWL API 4.5 · HermiT 1.4.3 · JavaFX 17 · Java 11 · Maven
}\\[1em]

\rule{\linewidth}{0.4pt}\\[0.4em]
{\normalsize Année universitaire 2025 -- 2026}

\end{center}
\end{titlepage}

% ============================================================
\section{Introduction}

Ce TP exploite plusieurs \emph{toolboxes} de représentation des connaissances
pour modéliser deux domaines réels et exécuter des requêtes de raisonnement
sur ces modèles.

\paragraph{Objectif.} Pour chaque formalisme, on doit (i)~choisir un domaine
réel, (ii)~représenter ses connaissances avec une syntaxe formelle adaptée,
(iii)~exploiter cette représentation pour répondre à des questions via un
moteur de raisonnement automatisé.

\paragraph{Périmètre.} Cinq formalismes ont été implémentés alors que le
minimum demandé en était deux. Chaque formalisme est interrogeable
interactivement depuis une interface graphique JavaFX unifiée
(\code{mvn javafx:run}) ou en console (\code{mvn exec:java}). Le code source
complet, ainsi qu'un rapport détaillé avec schémas TikZ, est disponible
séparément.

\paragraph{Choix des domaines.} Les deux domaines retenus pour traverser les
cinq formalismes sont :

\begin{itemize}
  \item \textbf{Industrie musicale algérienne} --- artistes, labels, albums,
        certifications, droits. Quatre rappeurs algériens (\emph{Soolking,
        Feu, Tif, Flenn}) servent d'individus communs à tous les formalismes,
        ainsi que deux labels (\emph{Universal} et \emph{Believe}) et deux
        albums (\emph{Guerilla} et \emph{Dernière~Heure}). Ce domaine met en
        évidence les notions de \emph{contrat}, \emph{indépendance}, et
        \emph{partage des droits} --- des situations naturellement non-monotones
        (un artiste indépendant qui signe un contrat).

  \item \textbf{Système judiciaire} --- acteurs (juges, avocats, procureurs,
        témoins), accusés et infractions. Trois accusés (\emph{Ali, Omar,
        Karim}) illustrent trois cas représentatifs : Ali (présumé innocent
        avec témoignage favorable), Omar (récidiviste, donc exception à la
        présomption d'innocence), et Karim (mineur poursuivi pour crime grave,
        situation où plusieurs règles s'opposent). Ce domaine est
        particulièrement adapté à la démonstration des exceptions et des
        conflits de règles.
\end{itemize}

\paragraph{Bibliothèques utilisées.} Voir la section
\emph{Bibliothèques et références} en fin de document pour la liste détaillée
avec liens.

\newpage
% ============================================================
\section{TP~1 --- Logique Modale}

\subsection*{Domaine choisi}

Univers musical à \textbf{quatre mondes} représentant des points de vue
distincts sur l'industrie de la musique : $w_1 = \emph{label}$
(maison de disque), $w_2 = \emph{fan}$, $w_3 = \emph{artiste}$, et
$w_4 = \emph{critique}$ musical. Chaque monde est caractérisé par les
propositions qui y sont vraies (par exemple : un album est platine, un
concert est annulé, un artiste est sous contrat).

\noindent La relation d'accessibilité $R$ traduit qui « voit »~quoi :
\begin{center}
\small $R = \{(w_1, w_2),\; (w_1, w_3),\; (w_3, w_3),\; (w_2, w_4),\; (w_4, w_4)\}$
\end{center}

\noindent Le label voit fans et artistes ; l'artiste se réfléchit lui-même
(boucle réflexive partielle) ; les fans accèdent aux critiques, etc. Ce
modèle illustre les notions de \emph{nécessité} (« dans tous les mondes
accessibles ») et de \emph{possibilité} (« dans au moins un »).

\subsection*{Représentation}

Modèle de Kripke $\mathcal{M} = \langle W, R, V\rangle$ implémenté avec
TweetyProject (\code{MlHerbrandInterpretation} pour les mondes,
\code{AccessibilityRelation} pour~$R$). Sémantique récursive :
\[
\mathcal{M}, w \models \nec\varphi \;\Leftrightarrow\; \forall v\;((w,v)\in R \Rightarrow \mathcal{M}, v \models \varphi)
\qquad
\mathcal{M}, w \models \pos\varphi \;\Leftrightarrow\; \exists v\;((w,v)\in R \land \mathcal{M}, v \models \varphi)
\]

\subsection*{Exploitation}

\begin{center}\small
\begin{tabular}{cllc}
\toprule
\textbf{\#} & \textbf{Formule} & \textbf{Monde} & \textbf{Résultat} \\
\midrule
Q1 & $\nec\,\text{album\_est\_platine}$       & $w_1$ & \textbf{FAUX} (faux en $w_2$) \\
Q2 & $\pos\,\text{concert\_annule}$           & $w_1$ & \textbf{VRAI} (vrai en $w_2$) \\
Q3 & $\nec\,\text{chanson\_en\_top\_charts}$ & $w_2$ & \textbf{VRAI} (vrai en $w_4$) \\
Q4 & $\pos\,\text{artiste\_independant}$     & $w_1$ & \textbf{VRAI} (vrai en $w_3$) \\
Q5 & $\nec(\text{album} \to \text{contrat})$ & $w_1$ & \textbf{FAUX} (faux en $w_3$) \\
\bottomrule
\end{tabular}
\end{center}

\noindent \textbf{Propriétés de $R$} (calculées automatiquement par la GUI) :
non réflexive, non transitive, non symétrique, mais \emph{sérielle}
(chaque monde a au moins un successeur).

\subsection*{Extension temporelle (cours chap. 3)}

La logique modale temporelle spécialise $R$ en \emph{ordre temporel} :
les mondes deviennent des \emph{instants}, et la relation $<$ correspond au
passage du temps. On ajoute quatre opérateurs :
\[
\textbf{G}\,p\ (\text{Going-to}) = \nec\,\text{futur}
\qquad
\textbf{F}\,p\ (\text{Futur}) = \pos\,\text{futur}
\qquad
\textbf{H}\,p\ (\text{Has-been}) = \nec\,\text{passé}
\qquad
\textbf{P}\,p\ (\text{Past}) = \pos\,\text{passé}
\]

\paragraph{Domaine musique --- carrière de Soolking.}
Cinq instants $t_1 < t_2 < t_3 < t_4 < t_5$ : 2017 (indépendant), 2018 (signe
chez Universal), 2019 (sort \emph{Guerilla}), 2021 (album platine), 2023
(nouvel album). Quelques requêtes :

\begin{center}\small
\begin{tabular}{cllc}
\toprule
\textbf{\#} & \textbf{Formule} & \textbf{Instant} & \textbf{Résultat} \\
\midrule
Q1 & $\mathbf{G}\,\text{platine}$            & $t_4$ & \textbf{VRAI} (vrai aussi en $t_5$) \\
Q2 & $\mathbf{F}\,\text{sous\_contrat}$      & $t_1$ & \textbf{VRAI} (signe Universal en $t_2$) \\
Q3 & $\mathbf{H}\,\text{artiste\_indep.}$    & $t_2$ & \textbf{VRAI} (vrai en $t_1$) \\
Q4 & $\mathbf{P}\,\text{album\_sorti}$       & $t_5$ & \textbf{VRAI} (sorti en $t_3$) \\
\bottomrule
\end{tabular}
\end{center}

\paragraph{Domaine justice --- déroulement d'un procès.}
Six instants $t_1 < \ldots < t_6$ : infraction, arrestation, inculpation,
audience, verdict, libération. Démonstration de \emph{changement de vérité
au cours du temps} :

\begin{center}\small
\begin{tabular}{cllc}
\toprule
\textbf{\#} & \textbf{Formule} & \textbf{Instant} & \textbf{Résultat} \\
\midrule
Q1 & $\mathbf{F}\,\text{condamne}$          & $t_1$ & \textbf{VRAI} (verdict en $t_5$) \\
Q2 & $\mathbf{H}\,\text{libre}$             & $t_3$ & \textbf{FAUX} (en détention dès $t_2$) \\
Q3 & $\mathbf{H}\,\text{presume\_innocent}$ & $t_4$ & \textbf{VRAI} (présomption d'innocence avant verdict) \\
Q4 & $\mathbf{F}\,\text{libre}$             & $t_3$ & \textbf{VRAI} (libéré en $t_6$) \\
\bottomrule
\end{tabular}
\end{center}

\noindent \emph{Implémentation :} sous-package \code{rcr.modal.temporal}
(\code{TemporalLogicMusic.java}, \code{TemporalLogicJustice.java}).
Évaluation directe en Java pur sur une liste ordonnée d'instants
(Tweety ne supporte pas nativement les opérateurs temporels).

\newpage
% ============================================================
\section{TP~2 --- Logique des Défauts}

\subsection*{Domaine choisi}

Industrie musicale algérienne, focalisée sur la question des droits d'auteur
et du statut contractuel des artistes. On modélise quatre artistes
(\emph{Soolking, Feu, Tif, Flenn}) avec leurs faits connus :

\begin{center}\small
\begin{tabular}{lll}
\toprule
\textbf{Artiste} & \textbf{Faits connus} & \textbf{Situation représentée} \\
\midrule
Soolking  & EstArtiste + SigneLabel + AlbumSorti       & Artiste signé avec album sorti \\
DZ\_Master & EstArtiste + IndepLongtemps + NouveauContrat & Indépendant qui vient de signer \\
Atlas     & EstArtiste + AlbumSorti + PlusDeDeuxAns    & Artiste expérimenté \\
\bottomrule
\end{tabular}
\end{center}

\subsection*{Représentation}

Théorie de défauts $\Delta = (W, D)$ où $W$ est l'ensemble des faits certains
et $D$ celui des règles par défaut. Chaque défaut s'écrit :
\[
\frac{\alpha \;:\; \beta}{\gamma}
\quad\text{(prérequis : justification / conclusion)}
\]

\noindent Les sept défauts du domaine :
\begin{align*}
d_1 &: \text{EstArtiste}(x) \,:\, \text{LibreDroits}(x) \,/\, \text{LibreDroits}(x) \\
d_2 &: \text{SigneLabel}(x) \,:\, \lnot\text{LibreDroits}(x) \,/\, \lnot\text{LibreDroits}(x) \\
d_3 &: \text{EstArtiste}(x) \land \text{AlbumSorti}(x) \,:\, \text{Royalties}(x) \,/\, \text{Royalties}(x) \\
d_4 &: \text{SigneLabel}(x) \,:\, \text{PartageRoyalties}(x) \,/\, \text{PartageRoyalties}(x) \\
d_5 &: \text{IndepLongtemps}(x) \,:\, \text{ExpSolo}(x) \,/\, \text{ExpSolo}(x) \\
d_6 &: \text{NouveauContrat}(x) \,:\, \lnot\text{LibreDroits}(x) \,/\, \lnot\text{LibreDroits}(x) \\
d_7 &: \text{PlusDe2Ans}(x) \land \text{AlbumSorti}(x) \,:\, \text{AlbumCertifie}(x) \,/\, \text{AlbumCertifie}(x)
\end{align*}

\noindent Implémentation via \code{RdlParser} et \code{SimpleDefaultReasoner}
de Tweety, qui calculent les \emph{extensions} (ensembles stables au sens de
Reiter).

\subsection*{Exploitation}

\textbf{Q1 --- Soolking} ($d_1$ et $d_2$ entrent en conflit) :
\textbf{2~extensions} stables sont calculées. L'une contient
\code{LibreDroits(soolking)}, l'autre $\lnot\text{LibreDroits(soolking)}$.
Les défauts $d_3$ et $d_4$ s'appliquent dans les deux $\Rightarrow$
\code{Royalties} et \code{PartageRoyalties} dans toutes les extensions.

\textbf{Q4 --- Atlas (démo non-monotonie).}
\begin{center}\small
\begin{tabular}{lll}
\toprule
\textbf{Théorie} & \textbf{Extensions} & \textbf{LibreDroits(atlas)} \\
\midrule
$W_0 = \{\text{EstArtiste, AlbumSorti, PlusDe2Ans}\}$ & 1 extension & \textbf{présent} ($\checkmark$) \\
$W_0 \cup \{\text{SigneLabel(atlas)}\}$               & 2 extensions & \textbf{plus dans toutes} ($\times$) \\
\bottomrule
\end{tabular}
\end{center}

\noindent L'ajout du fait \code{SigneLabel(atlas)} déclenche $d_2$ qui entre en
conflit avec $d_1$ : la conclusion \code{LibreDroits(atlas)}, valide dans la
première théorie, n'est plus dérivable dans toutes les extensions de la seconde.
C'est la \textbf{non-monotonie} du raisonnement par défaut.

\newpage
% ============================================================
\section{TP~3 --- Logique Classique (Propositionnelle \& Prédicats)}

\subsection*{Domaine choisi}

Le même quatuor d'artistes (\emph{Soolking, Feu, Tif, Flenn}) est modélisé
deux fois : (i)~en propositionnel, où chaque artiste est un atome
indépendant, (ii)~en prédicats, où l'on utilise un seul prédicat universel
$\forall x$ qui s'instancie sur les quatre individus.

\noindent Côté justice, trois accusés : \emph{Ali} (avec témoignage favorable),
\emph{Omar} (récidiviste), \emph{Karim} (mineur poursuivi pour crime grave).

\subsection*{Représentation propositionnelle (CNF + SAT4J)}

Faits (clauses unitaires) : \code{soolking}, \code{feu}, \code{tif},
\code{flenn}, \code{signe\_soolking}, \code{signe\_feu}, \code{indep\_feu}.
Clauses dérivationnelles :
\begin{align*}
c_1 &= \lnot\text{soolking} \lor \lnot\text{signe\_soolking} \lor \text{partage\_soolking} \\
c_2 &= \lnot\text{feu} \lor \lnot\text{signe\_feu} \lor \text{partage\_feu} \\
c_3 &= \lnot\text{tif} \lor \text{royalties\_tif} \\
c_4 &= \lnot\text{flenn} \lor \text{royalties\_flenn} \\
c_5 &= \lnot\text{indep\_feu} \lor \text{expsolo\_feu}
\end{align*}

\noindent Le test d'entailment $\text{KB} \models \varphi$ est réalisé par
\emph{réfutation} via \code{Sat4jSolver} (encapsulé dans \code{SatReasoner} de
Tweety) : $\text{KB} \models \varphi$ ssi $\text{KB} \cup \{\lnot\varphi\}$
est insatisfiable.

\subsection*{Représentation par prédicats (chaînage avant)}

Domaine fini $\{$soolking, feu, tif, flenn$\}$. Axiomes universels :
\begin{align*}
\text{ax}_1 &: \forall x\,(\text{EstRappeurAlgerien}(x) \Rightarrow \text{EstRappeur}(x)) \\
\text{ax}_2 &: \forall x\,(\text{EstRappeur}(x) \Rightarrow \text{EstArtiste}(x)) \\
\text{ax}_3 &: \forall x\,(\text{EstArtiste}(x) \Rightarrow \text{EstMusicien}(x)) \\
\text{ax}_5 &: \forall x\,(\text{SigneLabel}(x) \Rightarrow \text{PartageRoyalties}(x))
\end{align*}

\noindent Le chaînage avant instancie chaque axiome sur tous les individus et
dérive jusqu'au point fixe.

\subsection*{Exploitation}

\begin{center}\small
\begin{tabular}{ll}
\toprule
\textbf{Requête} & \textbf{Résultat} \\
\midrule
KB satisfiable ?                                       & \textbf{VRAI} \\
KB $\models$ partage\_soolking                         & \textbf{VRAI} (faits + $c_1$) \\
KB $\models$ expsolo\_feu                              & \textbf{VRAI} (faits + $c_5$) \\
KB $\cup\{\lnot\text{royalties\_flenn}\}$ satisfiable ? & \textbf{FAUX} (contradiction) \\
$\forall x\,\text{EstMusicien}(x)$                     & \textbf{VRAI} (chaîne ax$_1$+ax$_2$+ax$_3$) \\
$\exists x\,\text{ExperienceSolo}(x)$                  & \textbf{VRAI} (témoin : \emph{feu}) \\
\bottomrule
\end{tabular}
\end{center}

\noindent La GUI permet de tester en direct la satisfiabilité, l'entailment
$\text{KB} \models \varphi$ et l'ajout temporaire de formules pour observer les
contradictions provoquées.

\newpage
% ============================================================
\section{TP~4 --- Logique de Description (OWL DL + HermiT)}

\subsection*{Domaine choisi}

Ontologie musicale algérienne riche, avec une hiérarchie de concepts
(\code{Artiste} $\sq$ \code{Personne}, \code{Rappeur} $\sq$ \code{Artiste},
\code{RappeurAlgerien} $\sq$ \code{Rappeur}\dots), des rôles
(\code{signe\_chez}, \code{produit}, \code{recoit}\dots) et 10 individus
nommés :

\begin{center}\small
\begin{tabular}{lll}
\toprule
\textbf{Individu} & \textbf{Type asserté en ABox} & \textbf{Rôle} \\
\midrule
Soolking, Feu, Tif, Flenn   & RappeurAlgerien      & Rappeurs algériens \\
Universal                   & MajorLabel           & Major label international \\
Believe                     & IndependantLabel     & Label indépendant \\
Guerilla, Dernière Heure    & Album                & Albums produits \\
Cert\_Or, Cert\_Platine     & CertificationOr/Pl.  & Certifications de ventes \\
\bottomrule
\end{tabular}
\end{center}

\subsection*{Représentation (TBox)}

Quelques axiomes terminologiques caractéristiques (sur 25+ au total) :
\begin{align*}
\text{Rappeur}            &\equiv \text{Artiste} \sqcap \exists\,\text{pratique}.\text{Rap} \\
\text{RappeurAlgerien}    &\equiv \text{Rappeur} \sqcap \exists\,\text{chante\_en}.\text{Darija} \\
\text{ArtisteSigne}       &\equiv \text{Artiste} \sqcap \exists\,\text{signe\_chez}.\text{Label} \\
\text{ArtisteIndependant} &\equiv \text{Artiste} \sqcap \lnot\,\exists\,\text{signe\_chez}.\text{MajorLabel} \\
\text{AlbumOr}            &\equiv \text{AlbumCertifie} \sqcap \exists\,\text{recoit}.\text{CertificationOr}
\end{align*}

\noindent ABox typique : \code{signe\_chez(soolking, universal)},
\code{produit(soolking, guerilla)}, \code{recoit(guerilla, cert\_or)}, etc.
Construction via OWL API 4.5 ; raisonnement par tableaux avec HermiT.

\subsection*{Exploitation}

\begin{center}\small
\begin{tabularx}{\linewidth}{lcX}
\toprule
\textbf{Requête} & \textbf{Résultat} & \textbf{Justification (HermiT)} \\
\midrule
ArtisteIndependant(soolking)        & \textbf{FAUX} & signe\_chez universal:MajorLabel \\
AlbumOr(guerilla)                   & \textbf{VRAI} & recoit(guerilla, certOr:CertificationOr) \\
ArtisteMajor $\sq$ ArtisteSigne     & \textbf{VRAI} & MajorLabel $\sq$ Label \\
Disjoint(ArtisteIndep, ArtisteMajor) & \textbf{VRAI} & par construction des définitions \\
instances(ArtisteSigne)             & \{soolking, feu\} & signe\_chez vers un Label \\
\bottomrule
\end{tabularx}
\end{center}

\paragraph{Note OWA.} Sous l'\emph{Open World Assumption} (\textsc{owa}) d'OWL,
l'absence d'un fait n'est pas une négation. Pour les concepts définis par un
complément (ex. \code{ArtisteLibreDroits}), il a fallu insérer des assertions
explicites en ABox pour \emph{Tif} et \emph{Flenn}, équivalent à une fermeture
du monde locale.

\noindent La GUI calcule \emph{dynamiquement} la hiérarchie des classes via
\code{HermiT.getSubClasses} et l'affiche dans un \code{TreeView}. Trois modes
de requête interactive sont disponibles : Instance ($a:C$), Subsomption
($C \sq D$), et Disjonction ($C \sqcap D = \bot$).

\newpage
% ============================================================
\section{TP~5 --- Réseaux Sémantiques}

\subsection*{Domaine choisi}

Deux réseaux structurellement similaires illustrent la même intuition :

\begin{itemize}
  \item \textbf{Musique} : hiérarchie \code{Artiste $\to$ ArtisteSigne $\to$
        Soolking}, avec un \code{ArtisteIndep} qui hérite de \code{Artiste}
        mais possède une \emph{exception stricte} $\lnot\text{sous-contrat}$
        envers \code{Label}, annulant l'héritage général.
  \item \textbf{Justice} : hiérarchie \code{PartieProcès $\to$ Accusé $\to$
        Récidiviste}, où \code{Accusé} hérite de la \emph{présomption
        d'innocence} ; le \code{Récidiviste} la perd via une exception stricte
        $\lnot\text{presume}$.
\end{itemize}

\noindent Quatre styles d'arcs \code{is-a} (cours chapitre~5) :
\code{S} (strict, règle absolue), \code{NS} (non-strict, typique),
\code{ES} (exception stricte = lien barré), \code{EN} (exception non-stricte).
Les \emph{modalités épistémiques} (« le Procureur croit que\dots ») sont
construites structurellement par des nœuds-instances $C_1$, $N_1$ de classes
modales \code{Croyances}, \code{Negations}.

\subsection*{Représentation (extrait justice)}

\begin{center}\small
\begin{tabular}{lll}
\toprule
\textbf{Arc} & \textbf{Style} & \textbf{Sémantique} \\
\midrule
Accuse $\xrightarrow{\text{is-a}}$ PartieProces          & S    & Tout accusé est une partie au procès \\
Recidiviste $\xrightarrow{\text{is-a}}$ Accuse           & S    & Tout récidiviste est un accusé \\
Accuse $\xrightarrow{\text{presume}}$ Innocence          & rel. & Présomption d'innocence (héritée) \\
Recidiviste $\xrightarrow{\lnot\text{presume}}$ Innocence & ES   & \textbf{Exception stricte} : bloque l'héritage \\
\bottomrule
\end{tabular}
\end{center}

\subsection*{Algorithme d'héritage avec exceptions}

Notre moteur \code{inherits(n, r, t)} effectue un BFS depuis le nœud $n$ : à
chaque niveau, si un arc direct $\xrightarrow{r} t$ est trouvé, son signe
décide immédiatement (positif $\Rightarrow$ \code{INHERITED}, négatif
$\Rightarrow$ \code{BLOCKED}). Sinon on remonte les arcs \code{is-a}. C'est
le \textbf{principe de proximité} : la spécialisation la plus proche prime.

\begin{lstlisting}[caption={Moteur d'héritage (extrait)}]
public enum Inheritance { INHERITED, BLOCKED, NOT_FOUND }

public InheritanceTrace inherits(String n, String r, String t) {
    // BFS depuis n sur les arcs is-a
    // A chaque niveau, on cherche un arc direct (n', r, t)
    //   - positif -> INHERITED (decision proche prime)
    //   - negatif -> BLOCKED (exception ES/EN)
    // Sinon, on remonte. Si epuisement -> NOT_FOUND.
}
\end{lstlisting}

\subsection*{Exploitation}

\begin{center}\small
\begin{tabular}{lll}
\toprule
\textbf{Requête} & \textbf{Résultat} & \textbf{Pourquoi} \\
\midrule
Accuse presume Innocence ?      & \textbf{INHERITED} & arc direct sur Accuse \\
Recidiviste presume Innocence ? & \textbf{BLOCKED}   & exception $\lnot$presume au niveau Recidiviste \\
Soolking sous-contrat Label ?   & \textbf{INHERITED} & hérité via ArtisteSigne \\
ArtisteIndep sous-contrat Label ? & \textbf{BLOCKED} & exception stricte $\lnot$sous-contrat \\
\bottomrule
\end{tabular}
\end{center}

\noindent La GUI dessine le réseau avec un \emph{layout hiérarchique
auto-calculé}, des arcs en \emph{courbes de Bézier}, un \emph{hover} qui
illumine les arcs sortants, et colore les nœuds en vert (\code{INHERITED}) ou
en rouge (\code{BLOCKED}) lors d'une requête avec relation cible.

\newpage
% ============================================================
\section{Synthèse comparative}

\begin{center}\small
\begin{tabularx}{\linewidth}{lccccX}
\toprule
\textbf{Formalisme} & \textbf{Monotone} & \textbf{Hypoth.\ monde} & \textbf{Exceptions} & \textbf{Standard} & \textbf{Outil} \\
\midrule
Modale (Kripke)        & oui          & ---  & ---           & K, T, S4, S5      & Tweety MLogic \\
Défauts (Reiter)       & \textbf{non} & CWA  & multi-ext.    & Reiter 80         & Tweety RDL \\
Classique (Prop+FOL)   & oui          & CWA  & ---           & SAT, FOL          & SAT4J + Tweety FOL \\
Description (OWL DL)   & oui          & \textbf{OWA} & ---   & OWL 2 / W3C       & OWL API + HermiT \\
Réseaux sémantiques    & \textbf{non} & ---  & \textbf{ES/EN} & Quillian, Fahlman & Code Java ad hoc \\
\bottomrule
\end{tabularx}
\end{center}

\noindent Le cas \emph{Récidiviste} illustre la complémentarité : représenté
comme \emph{exception stricte} en réseaux sémantiques, comme \emph{défaut $d_6$
en conflit avec $d_1$} en logique des défauts, et comme \emph{classe
\code{AccuseRecidiviste} qui force \code{CasComplexe}} en logique de
description. Trois représentations différentes pour une même intuition
juridique.

\newpage
% ============================================================
\section{Bibliothèques et références}

\subsection*{Toolboxes et bibliothèques utilisées}

\begin{itemize}
  \item \textbf{TweetyProject 1.27}
        --- Toolbox Java pour les logiques computationnelles
        (modale, défauts, propositionnelle, FOL\dots).
        \href{https://tweetyproject.org/}{\nolinkurl{https://tweetyproject.org/}}\\
        Dépôt source :
        \href{https://github.com/TweetyProjectTeam/TweetyProject}%
             {\nolinkurl{https://github.com/TweetyProjectTeam/TweetyProject}}

  \item \textbf{OWL API 4.5.29}
        --- Bibliothèque de référence pour la création et la manipulation
        d'ontologies OWL~2 (W3C).
        \href{http://owlcs.github.io/owlapi/}{\nolinkurl{http://owlcs.github.io/owlapi/}}\\
        Dépôt source :
        \href{https://github.com/owlcs/owlapi}{\nolinkurl{https://github.com/owlcs/owlapi}}

  \item \textbf{HermiT 1.4.3.456}
        --- Raisonneur OWL~DL par tableaux, supportant la décidabilité de
        \textsc{sroiq}.
        \href{http://www.hermit-reasoner.com/}{\nolinkurl{http://www.hermit-reasoner.com/}}

  \item \textbf{SAT4J} --- Solveur SAT Java pur, intégré nativement à
        TweetyProject pour la propositionnelle.
        \href{https://www.sat4j.org/}{\nolinkurl{https://www.sat4j.org/}}

  \item \textbf{JavaFX 17} --- Framework graphique pour l'interface utilisateur.
        \href{https://openjfx.io/}{\nolinkurl{https://openjfx.io/}}

  \item \textbf{Maven 3.6+} --- Gestion de projet et dépendances.
        \href{https://maven.apache.org/}{\nolinkurl{https://maven.apache.org/}}

  \item \textbf{Java JDK 11} (OpenJDK).
        \href{https://openjdk.org/}{\nolinkurl{https://openjdk.org/}}
\end{itemize}

\subsection*{Documentation technique consultée}

\begin{itemize}
  \item Tweety --- documentation et tutoriels :
        \href{https://tweetyproject.org/doc/}{\nolinkurl{https://tweetyproject.org/doc/}}

  \item OWL 2 Web Ontology Language (W3C Recommendation) :
        \href{https://www.w3.org/TR/owl2-overview/}{\nolinkurl{https://www.w3.org/TR/owl2-overview/}}

  \item OWL API --- Javadoc et exemples :
        \href{https://owlcs.github.io/owlapi/apidocs_4/}%
             {\nolinkurl{https://owlcs.github.io/owlapi/apidocs_4/}}
\end{itemize}

\subsection*{Références théoriques}

\begin{itemize}
  \item \textbf{Reiter, R.} (1980). \emph{A Logic for Default Reasoning}.
        Artificial Intelligence, 13(1--2), 81--132.

  \item \textbf{Baader, F. et al.} (2003). \emph{The Description Logic Handbook}.
        Cambridge University Press.

  \item \textbf{Kripke, S.A.} (1963). \emph{Semantical considerations on modal
        logic}. Acta Philosophica Fennica, 16, 83--94.

  \item \textbf{Quillian, M.R.} (1968). \emph{Semantic Memory}.
        In : \emph{Semantic Information Processing}, MIT Press.

  \item \textbf{Fahlman, S.E.} (1979). \emph{NETL : a System for Representing
        and Using Real-World Knowledge}. MIT Press.

  \item Cours \textbf{RCR1} (Master 1 SII, USTHB, 2025--2026), notamment le
        chapitre~5 sur les réseaux sémantiques (conventions S / NS / ES / EN,
        construction structurelle des modalités).
\end{itemize}

\newpage
% ============================================================
\section{Conclusion}

Les cinq TPs montrent qu'aucun formalisme n'est universellement supérieur :
le choix dépend du domaine et des garanties souhaitées. La logique modale
capture les modalités (nécessité, possibilité) ; les défauts et les réseaux
sémantiques permettent la non-monotonie ; la logique classique offre une
sémantique simple mais rigide ; la logique de description hérite des
optimisations standardisées d'OWL~DL.

\medskip
\noindent Toutes les implémentations sont opérationnelles
(\code{mvn clean compile} $\to$ \textsc{build success}, 23 fichiers Java) et
interrogeables via la GUI unifiée (\code{mvn javafx:run}). Le code source
ainsi qu'un rapport détaillé sont disponibles dans le dépôt du projet.

\end{document}
